\chapter{Bevezető}
\label{chap:bevezeto}

% ── 1. RÉSZ: HÁTTÉR ÉS KONTEXTUS (1-2 bekezdés) ───────────────────────────────
% Indíts egy általános felvezetéssel a zenei információ-visszanyerésről (MIR) 
% és a hangszerfelismerés jelentőségéről a modern digitális világban.

Kodály Zoltán szerint a zene hangzó matematika. 
Hasonló analógiával élve, a számítástechnika pedig a digitális matematika egy formája, ami számomra azt mutatja meg, 
hogy valójában mindkét oldal közös tőről fakad.

Kutatásom és ezen dolgozat elkészítése pontosan ebben a metszéspontban, 
a zene és a számítástechnika határterületén valósul meg. 
Ezt a közös határterületet a tudományos világban a zenei információ-visszanyerés 
(Music Information Retrieval -- MIR) 
interdiszciplináris területe képviseli, 
amely a digitális jelfeldolgozás és a gépi tanulás eszköztárát ötvözi, 
hogy képessé tegye a számítógépes rendszereket az audiojelek automatikus 
értelmezésére. Ezen belül az automatikus hangszer- és hangszercsalád-felismerés az egyik legizgalmasabb kutatási irány, amely kulcsfontosságú a digitális zenei adatbázisok intelligens kereshetőségéhez, a valós idejű zeneelemzéshez, valamint az automatikus kottázáshoz (transzkripcióhoz).

Hétköznapi szinten ennek a technológiának az egyik leglátványosabb alkalmazása a modern zenei streaming szolgáltatók (mint például a Spotify vagy az Apple Music) ajánlórendszereiben érhető tetten. Ezek a platformok kifinomult algoritmusok segítségével elemzik és szabják személyre az általunk preferált zenei mintázatokat, hangszereket és dalokat, megteremtve ezzel a zenehallgatás új, intelligens korszakát. Ennek a személyre szabási folyamatnak a legalapvetőbb lépcsőfokát az olyan alacsony szintű zenei attribútumok és akusztikai jellemzők automatikus felismerése alkotja, mint amilyen a felvételeken megszólaló hangszerek azonosítása is.

A zenében a hangszerek változatos környezetben szólalhatnak meg. Háttér szempontjából beszélhetünk ,,laboratóriumi'' stúdiófelvételekről vagy valós idejű, környezeti háttérzajjal terhelt jelekről. A stúdióminőséghez jó analógia a légüres tér, ahol ideális körülmények között csak igencsak minimális zaj található meg, míg normális esetben szinte mindig zaj kíséri a zenét -- ha más nem, legalább olyan természeti hangok, mint a szél, a madárcsicsergés vagy a lépések zaja, illetve magának a hangszernek a fizikai megszólaltatása, ami a hangfeldolgozás esetében a második legnagyobb probléma. A hangszerek megszólalhatnak egyedül, de akár többen is egyszerre: több hangszer együttes megszólalása egymásra adódik, úgymond összeolvad, aminek a szétválasztása és megkülönböztetése -- mint a hangszerfelismerés ,,első számú károkozója'' -- különösen orchesztrák esetén szinte lehetetlen.

A probléma orvoslására az utóbbi időkben a deep learning modellek, különösen a 2D CNN hozott nagy könnyebbséget, amely a hang szeleteiből hoz létre egy-egy spektrális képet, mintegy a hangok ujjlenyomataként. A hangokat képpé alakítva, a vizuális mintázatok és karakterisztikák alapján történő elemzés lehetővé teszi a hangszerek sajátos színezetének felismerését, ám a szélsőségesen komplex zenekaroknál ez a módszer sem jelent teljes áttörést.

A kutatáshoz használt adathalmazt saját magam állítottam össze az interneten fellelhető zenefelvételek egészen apró szeleteiből, miután alaposan áttekintettem a korábban létező adatbázisokat. Azok a meglévő minták ugyanis nem követték pontosan az én tervezési szándékaimat, és nem adtak választ azokra a specifikus problémákra, amelyeket vizsgálni szerettem volna. Hosszú távon sokkal célravezetőbb egy teljesen saját, organikus adatbázis felépítése, amely tökéletesen illeszkedik a rendszer egyedi igényeihez.

Az elkészült neurális hálózatra és a saját adathalmazra építve a kutatás végcélja egy olyan alkalmazás kifejlesztése, amely közvetlen, gyakorlati segítséget nyújthat a felhasználóknak. Egy ilyen rendszer továbbfejlesztett változata különösen hasznos lehet azok számára, akik ritkán fordultak meg valós, élő zenei környezetben, és fejletlen a hangszerfelismerési képességük, de gyakran találják magukat olyan helyzetben, amikor egy számukra tetsző, videóklip nélküli dallamban megszólaló hangszer nevét szeretnék azonosítani.

A dolgozat további fejezetei lépésről lépésre követik a szoftverfejlesztés lépéseit. A zenei és gépi tanulási elmélet megalapozása után részletesen bemutatásra kerül a rendszer specifikációja és architektúrája, amelyet a 2D CNN hálózat és a valós idejű alkalmazás implementációja, valamint a gyakorlati tesztelés és kiértékelés követ.
